Early computational approaches to extra-classical receptive field (eCRF) effects primarily emphasized the role of intrinsic horizontal connections within the primary visual cortex (V1). These lateral projections allow for spatial integration by linking neurons with non-overlapping receptive fields across the cortical surface. However, horizontal connections face significant physiological constraints regarding their spatial extent and conduction velocity, which often fail to match the rapid, long-range dynamics of surround suppression observed experimentally. While lateral inhibition within V1 can produce a measurable decrease in firing rates, simulations often show that it typically yields a modest modulation—frequently failing to exceed a 30-35\% reduction in activity. This is insufficient to reach the robust levels of suppression (~75\%) reported in biological systems, suggesting that local V1 circuitry alone lacks the necessary dynamic range to fully account for global contextual integration.

While hierarchical feedback from higher-order areas such as V2 has been widely proposed as a driver of long-range contextual modulation, conventional computational implementations often rely on simplified feedback loops. These models typically treat top-down signals as linear additive inputs or non-specific gain modulators, which fail to capture the highly non-linear and stimulus-specific nature of extra-classical receptive field (eCRF) effects. Crucially, without internal competitive dynamics within the feedback-generating areas, top-down signals cannot achieve the precise spatial and feature-specific filtering required to match physiological observations. This limitation prevents existing architectures from reaching robust levels of surround suppression (~75\%), underscoring the necessity of a framework where hierarchical competition refines the feedback signals before they influence the laminar microcircuits of V1.
